\setcounter{section}{0}

\chapter*{Faire un plan}
\subsection*{\centering Objectif : Organiser son oral.}

\paragraph{Pourquoi faire un plan ? }
Le plan servira à structurer votre présentation en la découpant en temps et en idée.
Par rapport, à une rédaction complète, il présente l'avantage d'être un support qui évite la récitation et améliore la fluidité de la présentation.

\paragraph{Comment se structure un plan ? }
Un plan se découpe en grande partie (souvent introduction, plusieurs axes de développement et une conclusion).
À l'intérieur de ces parties, on trouve des points (\textit{bullets points} en anglais).

L'épreuve dure 20 min : 5 min de présentation de la question, 10 min d'échange avec le jury, puis 5 min sur le projet d'orientation.
Il est conseillé d'avoir une introduction, 2 axes de développement (3 au plus) et une conclusion pour organiser les 5 minutes de présentation.

\begin{center}
    La forme du plan conseillé est donc :
    \begin{enumerate}
        \item Introduction ($\approx 1min$)
        \item Axe 1 ($\approx 2min$)
        \item Axe 2 ($\approx 1min$)
        \item Conclusion ($\approx 1min$)
    \end{enumerate}    
\end{center}

\textbf{Un point est-un petit texte de 1 à 3 lignes maximums, par exemple : un argument, un pense-bête, une statistique à donner, une citation, une définition…
}.
Lorsque vous donnez une information, une statistique, définition, il faut en citer la source. 


\paragraph{Réflexion statistique :}
Lors d'une présentation orale, la vitesse d'élocution conseillée, c'est-à-dire le nombre de mots prononcés par minute pour une réception idéale, se situe autour des 100–125 mots par minute\footnote{Wong, Linda (2014). Essential Study Skills. Cengage Learning. ISBN 978-1285965628 \url{https://books.google.fr/books?id=yL2iAgAAQBAJ&q=conversation+words+per+minute&pg=PA348&redir_esc=y}.}.
Une phrase à l'oral contenant environ 7 mots, on considère donc \textbf{15-18 phrases par minutes}.
Ainsi, pour que votre plan soit un bon support oral, chaque ligne du plan doit pouvoir vous permettre de développer quelques phrases.

On peut tirer plusieurs conclusions de cette réflexion :
\begin{enumerate}
    \item \textbf{Chaque point doit être assez clair, concis et précis pour vous permettre de développer quelques phrases à son propos, disons entre 3 et 5 phrases (pertinentes)}.
    \item \textbf{Le plan devrait faire entre 15 et 30 points} — en fonction de votre aisance à l'oral\footnote{pour les curieux : $ \frac{15\ phrases \times 5\ minutes}{5\ phrases\ par\ point} = 15 points$ et $ \frac{18\ phrases \times 5\ minutes}{3\ phrases\ par\ point} = 30 points$} - pour couvrir les 5 minutes.
\end{enumerate}

\paragraph{Quelques astuces :}
\begin{itemize}
    \item Vous pouvez surligner d’une certaine couleur les mots clés, les arguments, les points forts de votre discours sur lesquels vous allez insister oralement.
    \item Vous pouvez surligner d’une autre couleur les éléments de type « enchaînement », « connecteurs », « annonce d'un exemple ou d’un argument », sur lesquels vous pouvez également parler un petit peu différemment ou prévoir des pauses puisque ce sont aussi des moments importants.
    \item Votre plan peut prendre différente forme, en outre du plan classique, vous pouvez par exemple utiliser une carte mentale ou un schéma.
    \item Les développements aussi peuvent être remplacés par des schémas, mais garder à l'esprit que vous n'aurez ni tableau noir ni diaporama : seules vos notes, avec éventuellement un schéma, sont autorisées.
    \item Vous devrez faire votre plan en 20 minutes de préparation de la présentation, il n'y a donc pas un plan, mais des plans et vous devez vous exercer à refaire votre plan plusieurs fois en vous chronométrant. 
\end{itemize}
\newpage

\section{L'introduction ($\approx 1min$)}

L'introduction est la première image que vous donnez de votre présentation c'est donc une partie primordiale.

\subsection*{Présentation succincte et rappel de la problématique}
Le but de ici est de vous présenter en 1 phrase et d'énoncer la problématique que vous allez traiter.
Contrairement au reste de votre présentation, je vous conseille de l'apprendre par cœur à l'avance. Il vous faut être sûr de vous lorsque vous énoncez votre problématique, \textbf{c'est votre sujet, vous ne pouvez pas le lire sur le papier}.

\subsection*{Motivations sur le choix du sujet}
La raison pour laquelle vous choisissez ce sujet en quelques phrases.

\subsection*{Accroche}
Le but de l'approche est d'essayer d’éveiller l’attention du jury par une accroche qui peut aussi vous permettre de montrer votre intérêt.
Suscitez la curiosité, donnez un exemple, un fait, une actualité, une image, une question, une citation, un lien personnel avec le sujet…

\subsection*{Contexte}
Le rappel du contexte est très dépendant du sujet :
\begin{description}
    \item[Si vous présentez un personnage :] Vous rappellerez brièvement ces attributs (date de naissance, date de mort, métier...) et le contexte historique. 
    \item[Si vous présentez une techonologie :] Vous rappellerez la période historique et les technologies utilisées avant sa création.
    \item[Si votre problématique contient des notions non-evidentes :] vous les définiraient en citant votre source.
\end{description}

\subsection*{Annonce du plan}
Vous annoncez votre plan en quelques phrases, vous pouvez en annoncer les grandes parties ou l’angle que vous avez choisi. Vous n’aurez pas le temps d’annoncer votre plan en détail.

\newpage

\section{Les axes de développements ($\approx 1$ à $2min$ chacun)}
Les axes de développements ont souvent une logique d'organisation commune qui permet leurs enchainements logiques.
Voici des exemples d'organisation des développements :
\begin{itemize}
    \item Thèse / antithèse / synthèse
    \item Avantages / inconvénients (ou l’inverse)
    \item Plan chronologique
    \item Constat / solutions
    \item Problèmes / causes / conséquences / solutions
    \item Auparavant / Aujourd’hui / Après
    \item Enchaînement logique d’idées, arguments et exemples
    \item Du général au particulier ou du particulier au général…
\end{itemize}

\subsection*{Argument}
Un des objectifs du Grand oral est de montrer votre « force de conviction ».
Il existe plusieurs sortes d’arguments que je conseille de panacher au sein de votre exposé (et de la discussion). Certains sont difficilement contestables, d’autres peuvent être contredits.
Dans tous les cas, il faut indiquer clairement la provenance de votre argument, le \textbf{sourcer}.
\begin{description}
    \item[Les arguments de droit] qui font référence à la loi.
    \item[Les arguments de fait] qui évoquent une réalité connue de tous, un fait,
un évènement, une étude, le résultat d’une expérience ou d’une observation…
    \item[Les arguments scientifiques] qui font consensus auprès de la communauté
scientifique et qui sont décrits par des lois ou par des théorèmes, par exemple.
    \item[Les arguments logiques] qui font appel à la raison ; les arguments de valeur qui invoquent les grands principes ou les fondements de notre société (morale, liberté…)
    \item[Les arguments d’autorité] qui s'appuient sur une personne reconnue ou une « Institution » « Descartes pense que… », « Sur le site de l’INSEE… »)
    \item[les arguments tirés de votre expérience personnelle ou de celle d’autrui] dans ce cas-là, dans la plupart des cas, on s'exprime sur son ressenti ou sur son opinion.
\end{description}

\section{La conclusion ($\approx 1min$)}

Selon les études, les cerveaux humains ont tendance à enregistrer une expérience en grande partie sur la base de ce qu'ils ont ressenti à son apogée (c'est-à-dire son point le plus intense) et à sa fin.
L’effet se produit indépendamment du fait que l’expérience soit agréable ou désagréable\footnote{Fredrickson, Barbara L.; Kahneman, Daniel (1993). "Duration neglect in retrospective evaluations of affective episodes". Journal of Personality and Social Psychology. 65 (1): 45–55. doi:10.1037/0022-3514.65.1.45. PMID 8355141.\url{https://psycnet.apa.org/doiLanding?doi=10.1037\%2F0022-3514.65.1.45}}.
Annoncez clairement votre conclusion : signalez au jury que vous avez quasiment terminé. Dites, par exemple : \og pour terminer… \fg, \og pour conclure… \fg, \og enfin… \fg, \og en conclusion \fg.

\subsection*{Réponse à la problématique}
Répondez à nouveau durant la conclusion, si vous avez posé une question qui attend une réponse (« pour ou contre… »).
Vous pouvez reprendre les termes de la question pour synthétiser la réponse que vous avez apportée.

\subsection*{Lien avec le projet de l'élève [Facultatif]}
Ici, vous pouvez faire le lien avec votre projet d'orientation si cela est pertinent ou discuter en quoi préparer cette présentation vous a apporté.

\subsection*{Ouverture}
Suggérez des idées, des pistes, proposez une ouverture, posez une autre question, avec un autre point de programme d’une de vos spécialités.
Cela vous permet d'anticiper les questions du jury est donc d'être prêt.e à l'avance.

